\documentclass[12pt]{article}
\usepackage{geometry}
\geometry{a4paper}
\usepackage{fancyhdr}
\usepackage{graphicx}
\usepackage{wallpaper}
\usepackage{Sweave}
\begin{document}

\pagestyle{fancy}
\ThisULCornerWallPaper{1.0}{./verticalbar.eps} % adds vertical bar for title page
\URCornerWallPaper{0.25}{./ctb_logo_uright.eps} %adds contrabass logo for title page - only suitable for bottom right - not appropriate for top right
\begin{titlepage}
\thispagestyle{fancy}
\fancyhead{}
\renewcommand{\headrulewidth}{0.0pt}
\fancyfoot{}
\lfoot{\footnotesize Last Revision: \today}
\rfoot{\footnotesize confidential}
\begin{center} % client logo at the center top area
\includegraphics[width=0.3\textwidth]{./AVS_logo.eps}\\[1cm]
\textsc{\LARGE Anvis Group}\\[1.5cm]
\textsc{\Large Transfer Pricing Policy Review}\\[0.5cm]
\huge\bfseries Outside Evidence Work\\[0.8cm]
\end{center}
\begin{abstract}
Anvis Group's transfer pricing policy generally complies with the OECD Transfer Pricing Guidelines (2010). In order to effectively limit the administrative burden associated with identifying and maintaining the required outside evidence a transactional net margin approach is chosen. Ultimately,the arm's length operating profit is employed to arrive at appropriate gross profit margins.\\[0.2cm] Several TP Catalyst studies (Bureau von Dijk) were conducted to identify contract services and contract R\&D comparables. The present paper documents the econometric analyses conducted to arrive at the final arm's length ranges used for transfer pricing policy purposes.\\[1cm] \end{abstract}
\begin{minipage}[b]{0.4\textwidth}
\begin{flushleft}\large
\emph{Author:}\\
Dr. Steffen R. \textsc{M\"ockel}\\[0.15cm]
\small Diplom-Volkswirt / Economist
\end{flushleft}
\end{minipage}
\hfill
\begin{minipage}[t]{0.5\textwidth}
\begin{flushleft}
\small
\emph{Contact:}\\
Konkordiastrasse 2 \\40129 D\"usseldorf\\[0.25cm]Tel.: +49 211 9304209\\ Mob.: +49 15774 111022\\[0.25cm]e-mail: dr.s.moeckel@contrabass.biz\\www.contrabass.biz
\end{flushleft}
\end{minipage}
\vfill
\end{titlepage}
\cfoot{\footnotesize confidential}
\rfoot{\footnotesize \thepage}
\section{Final arm's length ranges and their application}
\subsection{General provisions}
Anvis Group's transfer pricing system for intra-group services and contract R\&D relies on the transactional net margin method. At budget, the arm's length transfer price for these items is determined using a profit mark-up consistent with the mid-point of the respective interquartile range. The actual result is regularly reviewed (at least after 6 months). A transfer pricing adjustment is made whenever the respective service entity's actual profit is not within the corresponding arm's length range. The adjustment is made to the range's upper or lower bound whatever is closest to the actual result achieved.\\[0.2cm]
The adjustment is determined as retro-active adjustment ensuring that the results actually achieved are within the arm's length range. This is largely consistent with Anvis Group's current practice. In order to charge each recipient with the actual cost incurred on its behalf/the actual benefits received, Anvis Deutschland GmbH (AVS002) issues already debit or credit notes.\\[0.2cm]
As outlined below, the arm's length ranges reflect period averages. Accordingly, a transfer pricing adjustment back to the closest upper or lower bound can be based on period averages as well. In order to limit the transfer pricing related worklaod the comparison will be based the respective Anvis entities' annual data.\\[0.2cm]
The arm's length ranges are to be reviewed regularly.\begin{itemize}
  \item just adding new data (effectively employ a moving average approach)
  \item conduct a completely new comparables search
  \item recalculate the necessary comparability adjustments
\end{itemize}
In order to limit the transfer pricing related administrative burden it may be suitable to largely integrate all necessary functional adjustments into TP Catalyst. Additional statistical analysis should be conducted only if w/o such more advance analysis the arm's length functional profits are not in line with the Anvis Group's overall profit potential.
\section{Resolving the comparables' search challenges}
\subsection{The NACE based search strategy is supplemented}
\subsection{A number of contract service providers covers the services rendered}
\subsection{Contract R\&D comparables are taken from other industries}
\subsection{The evidence available takes the supplier's position}
\section{The contract service comparables}
\section{The Contract R\&D comparables}
The data represent a combination of panel and timeseries data. 
\begin{figure}
\includegraphics{AVS_OS_Docu-002}
\caption {Test}
\end{figure}
\end{document}
